\documentclass[11pt]{article}
\usepackage[margin=1in]{geometry}
\usepackage{enumitem}
\setlength{\emergencystretch}{3em}
% =============================================================================
% Preambles/header.tex — shared LaTeX/Beamer preamble for this project
%
% Use from a Beamer lecture by prepending:
%     \input{header}
% and compiling with TEXINPUTS=../Preambles:$TEXINPUTS (the /compile-latex
% skill does this automatically).
%
% PALETTE CONTRACT. Color names defined here MUST match the names in
% Quarto/theme-template.scss so Beamer and Quarto renderings use the same
% palette. The scripts/check-palette-sync.sh script enforces this.
%
% To customize: edit the HEX values below AND the matching $variable in
% Quarto/theme-template.scss, then run scripts/check-palette-sync.sh.
% =============================================================================

% -----------------------------------------------------------------------------
% Engine detection + encoding
%
% Our /compile-latex skill uses XeLaTeX, where inputenc is unnecessary and
% can produce encoding warnings. Only load inputenc under pdfTeX.
% -----------------------------------------------------------------------------
\usepackage{iftex}
\ifPDFTeX
  \usepackage[utf8]{inputenc}
\fi

\usepackage{amsmath, amssymb, amsthm}
\usepackage{booktabs}
\usepackage{graphicx}

% -----------------------------------------------------------------------------
% PALETTE — mirrors Quarto/theme-template.scss variable names.
% Load xcolor and define colors BEFORE anything that references them
% (notably hyperref, hypersetup, and the Beamer color assignments).
% -----------------------------------------------------------------------------
\usepackage{xcolor}

\definecolor{primary-blue}{HTML}{012169}      % headings, accents
\definecolor{primary-gold}{HTML}{B9975B}      % emphasis, borders
\definecolor{highlight-yellow}{HTML}{F2A900}  % markers, alerts
\definecolor{light-bg}{HTML}{E8EDF5}          % title / section backgrounds
\definecolor{jet}{HTML}{1A1A1A}               % body text

% Semantic colors (used in diagrams and annotations)
\definecolor{positive}{HTML}{15803D}          % good / observed / identified
\definecolor{negative}{HTML}{B91C1C}          % bad / problematic / bias
\definecolor{neutral}{HTML}{525252}           % reference / context

% Highlight accents (match .hi-* classes in the SCSS)
\definecolor{hi-slate}{HTML}{314F4F}
\definecolor{hi-green}{HTML}{2E8B57}
\definecolor{hi-red}{HTML}{C41E3A}

% Semantic aliases so skill templates and snippets can write
% \textcolor{accent}{...} without hard-coding "primary-blue".
\colorlet{accent}{primary-blue}
\colorlet{accent2}{primary-gold}

% -----------------------------------------------------------------------------
% Hyperref — loaded AFTER xcolor + palette so link colors resolve.
% -----------------------------------------------------------------------------
\usepackage{hyperref}
\hypersetup{colorlinks=true, linkcolor=primary-blue, urlcolor=primary-blue,
            citecolor=primary-blue, pdfborder={0 0 0}}

% -----------------------------------------------------------------------------
% Course code — set once per deck (after \input{header}, before \begin{document})
% via \coursecode{CS401}. Shown in the footer on every slide so a deck's course
% is identifiable without opening the title page. Empty by default.
% -----------------------------------------------------------------------------
\makeatletter
\newcommand{\coursecode}[1]{\gdef\@coursecodetext{#1}}
\gdef\@coursecodetext{}
\makeatother

% -----------------------------------------------------------------------------
% Beamer theme assignments (only apply under Beamer).
%
% We detect Beamer via \@ifclassloaded{beamer}, which is the documented,
% non-internal way. This must be wrapped in \makeatletter/\makeatother
% because \@ifclassloaded contains an @-catcode character.
% -----------------------------------------------------------------------------
\makeatletter
\@ifclassloaded{beamer}{%
  \usefonttheme{professionalfonts}
  \setbeamercolor{structure}{fg=primary-blue}
  \setbeamercolor{frametitle}{fg=primary-blue}
  \setbeamercolor{title}{fg=primary-blue}
  \setbeamercolor{subtitle}{fg=primary-gold}
  \setbeamercolor{author}{fg=primary-blue}
  \setbeamercolor{institute}{fg=neutral}
  \setbeamercolor{date}{fg=primary-gold}
  \setbeamercolor{itemize item}{fg=primary-blue}
  \setbeamercolor{itemize subitem}{fg=primary-gold}
  \setbeamercolor{alerted text}{fg=primary-gold}
  \setbeamercolor{block title}{fg=white, bg=primary-blue}
  \setbeamercolor{block body}{bg=light-bg}
  % Semantic block variants: block=definition/concept (blue), exampleblock=
  % worked example (green/positive), alertblock=key takeaway or caution (gold).
  % Keep to <=2 colored boxes per slide across ALL three variants (INV-7).
  \setbeamercolor{block title example}{fg=white, bg=positive}
  \setbeamercolor{block body example}{bg=light-bg}
  \setbeamercolor{block title alerted}{fg=white, bg=primary-gold}
  \setbeamercolor{block body alerted}{bg=light-bg}

  % Minimal footer: course code (left) + page X / N (right), no navigation clutter
  \setbeamertemplate{navigation symbols}{}
  \setbeamertemplate{footline}{%
    \color{neutral}\footnotesize\hspace{0.7em}\@coursecodetext\hfill
    \insertframenumber/\inserttotalframenumber\hspace{0.7em}\vspace{0.3em}}
}{}
\makeatother

% -----------------------------------------------------------------------------
% TikZ setup — libraries the snippet gallery uses
% -----------------------------------------------------------------------------
\usepackage{tikz}
\usetikzlibrary{arrows.meta, positioning, calc, decorations.pathreplacing,
                fit, shapes.geometric, backgrounds}

% Shared TikZ styles — snippets and hand-written diagrams can pick these up
% via `\begin{tikzpicture}[every node/.style={...}]` or by naming specific
% styles (e.g., `\node[dag-node] (X) at (0,0) {X};`).
\tikzset{
  % Node styles for causal DAGs and similar
  dag-node/.style = {
    draw=primary-blue, fill=light-bg, rounded corners=2pt,
    minimum width=1.2cm, minimum height=0.9cm, align=center, font=\small
  },
  decision-node/.style = {
    draw=primary-gold, fill=white, diamond, aspect=1.8,
    minimum width=1.8cm, minimum height=1.2cm, align=center, font=\small,
    inner sep=1pt
  },
  % Edge styles — observed vs counterfactual
  observed-edge/.style = {-{Latex[length=2.5mm]}, thick, draw=jet},
  counterfactual-edge/.style = {-{Latex[length=2.5mm]}, thick, draw=neutral, dashed},
  confound-edge/.style = {-{Latex[length=2.5mm]}, thick, draw=negative, dashed},
  % Data-point styles for plots
  observed-dot/.style = {fill=primary-blue, draw=primary-blue, circle, inner sep=1.2pt},
  counterfactual-dot/.style = {fill=white, draw=neutral, circle, inner sep=1.2pt}
}

% -----------------------------------------------------------------------------
% Convenience macros
% -----------------------------------------------------------------------------
\newcommand{\muted}[1]{\textcolor{neutral}{#1}}
\newcommand{\key}[1]{\textcolor{primary-gold}{\textbf{#1}}}
\newcommand{\good}[1]{\textcolor{positive}{#1}}
\newcommand{\bad}[1]{\textcolor{negative}{#1}}

% Standout frame macro: full-bleed dark background for section transitions.
% Beamer-only (uses \begin{frame}/\setbeamercolor/\usebeamerfont) -- do not
% call from an article-class file (e.g. Notes/<CODE>/*-notes.tex). \muted,
% \key, \good, \bad, and \coursecode above are plain xcolor/gdef and are
% safe to use from either class; verified when Notes/article-class support
% was added.
% Expand: \transitionslide{Your transition title}
\newcommand{\transitionslide}[1]{%
  \begingroup
  \setbeamercolor{background canvas}{bg=primary-blue}
  \begin{frame}[plain]
    \centering
    \vfill
    {\usebeamerfont{title}\color{white}\Large #1\par}
    \vfill
  \end{frame}
  \endgroup
}

% Section divider frame (Beamer-only), an alternative to \transitionslide with a
% darker two-line style: near-black (jet) background, a small white label line
% above a larger gold title, and a thin gold rule along the bottom edge.
% Expand: \sectiondivider{Section 2}{Pointers}
\newcommand{\sectiondivider}[2]{%
  \begingroup
  \setbeamercolor{background canvas}{bg=jet}
  \begin{frame}[plain]
    \vspace*{3.0cm}
    {\color{white}\ttfamily\large #1\par}
    \vspace{0.5cm}
    {\color{primary-gold}\LARGE #2\par}
    \begin{tikzpicture}[remember picture, overlay]
      \draw[primary-gold, line width=2.5pt]
        ($(current page.south west)+(0,0.1)$) -- ($(current page.south east)+(0,0.1)$);
    \end{tikzpicture}
  \end{frame}
  \endgroup
}

% -----------------------------------------------------------------------------
% Channel watermark — "Concepts That Click" (YouTube).
%
% Article-class files (CompetitiveExam Q/A sets, Notes, Assignments) get a
% light diagonal draft watermark on every page plus a centered footer naming
% the channel. Beamer decks get a subtle TikZ background watermark instead
% (the footline already carries the course code). Centralized here so every
% MCQ PDF is branded without per-file edits.
% -----------------------------------------------------------------------------
\makeatletter
\@ifclassloaded{article}{%
  \usepackage{draftwatermark}
  \SetWatermarkText{Concepts That Click}
  \SetWatermarkScale{1}
  \SetWatermarkFontSize{2cm}
  \SetWatermarkLightness{0.86}
  \SetWatermarkAngle{45}
  \usepackage{fancyhdr}
  \pagestyle{fancy}
  \fancyhf{}
  \fancyfoot[C]{\footnotesize\textcolor{neutral}{Concepts That Click~|~YouTube: \href{https://www.youtube.com/@concepts.thatclick}{@concepts.thatclick}}}
  \renewcommand{\headrulewidth}{0pt}
  \renewcommand{\footrulewidth}{0pt}
  \fancypagestyle{plain}{%
    \fancyhf{}%
    \fancyfoot[C]{\footnotesize\textcolor{neutral}{Concepts That Click~|~YouTube: \href{https://www.youtube.com/@concepts.thatclick}{@concepts.thatclick}}}%
    \renewcommand{\headrulewidth}{0pt}%
    \renewcommand{\footrulewidth}{0pt}%
  }
}{}
\@ifclassloaded{beamer}{%
  \setbeamertemplate{background}{%
    \begin{tikzpicture}[remember picture, overlay]
      \node[rotate=45, opacity=0.055, align=center]
        at (current page.center)
        {\Huge\textcolor{neutral}{Concepts That Click}};
    \end{tikzpicture}%
  }
}{}
\makeatother 
\coursecode{JKSSB}
\renewcommand{\thesection}{26.\arabic{section}}
\newtheorem{example}{Example}[section]
\newenvironment{definitionbox}[1]{%
  \par\noindent\textbf{Definition (#1).}\ \itshape
}{\par\vspace{0.3em}}
\title{Excel Functions --- Detailed Lecture Notes}
\author{JKSSB Computer Awareness}
\date{\today}

\begin{document}
\maketitle

\section{From formulas to functions}
MS Excel is built around the idea that a value can be \emph{calculated}
instead of typed. A cell that begins with the equal sign \texttt{=} holds a
calculation, and Excel works out the result. Such a calculation is called a
\textbf{formula}. A simple formula is \texttt{=A1+A2}, which adds the values
of two cells. A formula can also use operators such as \texttt{+}, \texttt{-},
\texttt{*}, and \texttt{/} to add, subtract, multiply, and divide.

Typing every cell into a formula becomes slow and error-prone when there are
many cells. Excel therefore provides \textbf{functions}: ready-made, named
operations that do a common job for you. \texttt{SUM} adds a range of
numbers, \texttt{AVERAGE} finds their mean, and \texttt{MAX} finds the largest.
A function is always written \emph{inside} a formula. The function name is
followed by parentheses that contain the \emph{arguments} --- the cells or
values the function should work on.

\begin{definitionbox}{Formula}
An expression that begins with the equal sign \texttt{=} and calculates a
value, for example \texttt{=A1+A2} or \texttt{=SUM(A1:A10)}.
\end{definitionbox}

\begin{definitionbox}{Function}
A predefined, named built-in operation, such as \texttt{SUM}, \texttt{COUNT},
or \texttt{IF}, that is used inside a formula to perform a particular
calculation on its arguments.
\end{definitionbox}

The distinction to hold is simple. A \textbf{formula} is the whole
calculation. A \textbf{function} is a built-in tool that performs one part of
that calculation. Every function is written as a formula, but not every
formula contains a function: \texttt{=A1+A2} is a formula with no function in
it, while \texttt{=SUM(A1:A10)} is a formula that uses one function. A
formula must begin with \texttt{=}; if the equal sign is missing, Excel treats
the entry as plain text and shows it as typed.

\section{Function syntax and arguments}
Every function follows the same shape: a name, then a pair of parentheses, and
inside the parentheses the arguments separated by commas. This is written as
\texttt{=FUNCTION(argument1, argument2)}. The \textbf{name} tells Excel which
operation to perform. The \textbf{arguments} tell it which cells or values to
use. For example, in \texttt{=SUM(A1:A10)} the name is \texttt{SUM} and the
single argument is the range \texttt{A1:A10}.

A \textbf{range} is a block of cells written with a colon between the first
and last cell. \texttt{A1:A10} means the cells A1, A2, A3, and so on down to
A10. A function can take more than one argument: \texttt{=SUM(A1:A5, C1:C5)}
adds the two ranges together. The comma separates one argument from the next.
Knowing the name, the parentheses, and the colon is enough to read any of the
basic functions in this lesson.

\begin{center}
\begin{tabular}{p{0.15\textwidth}p{0.28\textwidth}p{0.47\textwidth}}
\toprule
\textbf{Function} & \textbf{Example} & \textbf{What it returns} \\
\midrule
\texttt{SUM} & \texttt{=SUM(A1:A10)} & The total of the numbers in the range. \\
\texttt{AVERAGE} & \texttt{=AVERAGE(A1:A10)} & The arithmetic mean of the numbers. \\
\texttt{COUNT} & \texttt{=COUNT(A1:A10)} & The number of cells holding numbers. \\
\texttt{COUNTA} & \texttt{=COUNTA(A1:A10)} & The number of non-empty cells. \\
\texttt{MAX} & \texttt{=MAX(A1:A10)} & The largest number in the range. \\
\texttt{MIN} & \texttt{=MIN(A1:A10)} & The smallest number in the range. \\
\texttt{IF} & \texttt{=IF(A1>50,"Y","N")} & One of two values, chosen by the test. \\
\bottomrule
\end{tabular}
\end{center}

\section{SUM and AVERAGE}
\texttt{SUM} returns the total of the numbers in its arguments. Its most
common form is \texttt{=SUM(A1:A10)}, which adds all the numbers in the range
A1:A10. Cells that hold text and cells that are blank are ignored; only
numbers are added. SUM accepts several ranges, so \texttt{=SUM(A1:A5,B1:B5)}
adds both blocks in one step.

\texttt{AVERAGE} returns the arithmetic mean of the numbers in its arguments.
\texttt{=AVERAGE(A1:A10)} adds the numbers and divides by how many numbers
there are. The important detail is what counts in the divisor: blank cells and
cells holding text are \emph{not} counted. AVERAGE divides by the number of
\emph{numeric} cells, not by the total number of cells in the range. This is
the source of many wrong answers.

\begin{example}
Suppose A1:A5 holds the values 10, 20, and 30, while A4 and A5 are blank.
\texttt{=SUM(A1:A5)} returns $60$. \texttt{=AVERAGE(A1:A5)} returns
$(10+20+30)/3 = 20$, not $60/5 = 12$. Excel divides by the three cells that
actually hold numbers, because the two blanks are ignored.
\end{example}

\section{COUNT and COUNTA: the near-neighbour pair}
\texttt{COUNT} and \texttt{COUNTA} look almost the same but count different
things, and the difference is a favourite exam trap. \texttt{COUNT(counts)}
only the cells in a range that contain \emph{numbers}. \texttt{COUNTA} (the
``A'' stands for \emph{all}) counts every cell that is \emph{not empty},
whether it holds a number or text. Both functions ignore blank cells.

\begin{definitionbox}{COUNT}
A function that returns the number of cells in a range that contain numeric
values. Text, blanks, and logical values are not counted.
\end{definitionbox}

\begin{definitionbox}{COUNTA}
A function that returns the number of cells in a range that are not empty. It
counts numbers, text, and error values alike; only truly blank cells are
ignored.
\end{definitionbox}

The contrast is easiest to see with a mixed range. Suppose A1:A10 holds six
numbers, two text labels, and two blank cells. \texttt{=COUNT(A1:A10)}
returns 6, because only the numeric cells are counted. \texttt{=COUNTA(A1:A10)}
returns 8, because the six numbers and the two text labels are all non-empty.
Both functions ignore the two blanks. A related function, \texttt{COUNTBLANK},
does the opposite and counts the empty cells, returning 2 for the same range.
For the exam, the pair to memorise is: \textbf{COUNT = numbers only; COUNTA =
all non-empty cells}.

\section{MAX and MIN}
\texttt{MAX} returns the largest number in a range, and \texttt{MIN} returns
the smallest. Their forms are \texttt{=MAX(C1:C10)} and, for the smallest,
\texttt{=MIN(C1:C10)}. Like SUM, AVERAGE, and COUNT, they work on numbers and
ignore text and blank cells. A useful way to remember them is that \texttt{MAX}
is the maximum and \texttt{MIN} is the minimum; do not confuse \texttt{MIN}
with the word ``minimum cell'', which is not a function at all.

\begin{example}
A range holds the marks 42, 67, 55, 91, and 38. \texttt{=MAX(C1:C5)} returns
$91$, the top mark. \texttt{=MIN(C1:C5)} returns $38$, the lowest mark.
\texttt{=AVERAGE(C1:C5)} returns $58.6$, the mean of the five numbers.
\end{example}

\section{The IF function and nested IF}
\texttt{IF} lets a cell show one of two results depending on whether a
condition is true or false. Its syntax is
\texttt{=IF(logical\_test, value\_if\_true, value\_if\_false)}. The first
argument is a \emph{logical test} that can be answered only with true or
false, such as \texttt{A1>50}. The second argument is the value or text to
return when the test is true, and the third is the value to return when it is
false. Thus \texttt{=IF(A1>50,"Pass","Fail")} shows \texttt{Pass} when A1 is
more than 50 and \texttt{Fail} otherwise.

A \textbf{nested IF} places a second IF inside the value-if-false part of the
first, so more than one condition can be tested. A typical form is
\texttt{=IF(A1>50,"High",IF(A1>25,"Medium","Low"))}. Excel checks the
conditions in the order they are written, and the first test that is true
decides the result.

\begin{example}
Take the formula \texttt{=IF(A1>50,"High",IF(A1>25,"Medium","Low"))} and let
A1 be blank. In a numeric comparison Excel reads a blank cell as $0$. Since
$0$ is neither greater than 50 nor greater than 25, both tests are false and
the formula returns the final value, \texttt{Low}. This is the behaviour the
Junior Assistant 2026 paper examined (KB PYQ-14, JA 2026 Q74).
\end{example}

When the input is \emph{non-numeric} text, the behaviour depends on what the
formula does with it. A comparison of text against a number is still answered
by Excel, but if the formula tries to \emph{calculate} with text --- for
example, \texttt{=A1*2} where A1 holds a word --- Excel cannot do the
arithmetic and returns the error value \texttt{\#VALUE!}. The safe exam
reading is: a blank is treated as 0 in a numeric test, and text that is
arithmetic-only produces \texttt{\#VALUE!}.

\section{A worked sheet, function by function}
A small sheet makes the behaviour of every function concrete. Column A holds
five student names --- Anita, Bilal, Chandan, Divya, and Esha --- and column B
holds the marks 40, 55, 70, and 60, with Esha's mark (B5) left blank because
she has not been marked yet. Reading the sheet function by function shows what
each one does with the blank.

\begin{center}
\begin{tabular}{p{0.30\textwidth}p{0.30\textwidth}p{0.28\textwidth}}
\toprule
\textbf{Formula} & \textbf{Result} & \textbf{Why} \\
\midrule
\texttt{=SUM(B1:B5)} & 225 & Adds 40, 55, 70, 60; the blank is ignored. \\
\texttt{=AVERAGE(B1:B5)} & 56.25 & Divides 225 by 4 numeric cells, not by 5. \\
\texttt{=COUNT(B1:B5)} & 4 & Counts the four cells that hold numbers. \\
\texttt{=COUNTA(B1:B5)} & 4 & Counts the four non-empty marks; B5 is blank. \\
\texttt{=MAX(B1:B5)} & 70 & The largest of the four marks. \\
\texttt{=MIN(B1:B5)} & 40 & The smallest of the four marks. \\
\texttt{=IF(B5>=40,"Pass","Fail")} & Fail & B5 is blank, read as 0, so the test fails. \\
\bottomrule
\end{tabular}
\end{center}

The last row links the sheet back to the nested-IF behaviour: a blank cell is
read as zero in a numeric test, so Esha is shown as not yet passing. Note also
that \texttt{COUNTA} on the names column, \texttt{=COUNTA(A1:A5)}, would
return 5, because all five names are present even though only four marks are.
This is exactly how the two count functions are used together to spot missing
data.

\section{Other basic functions worth knowing}
Beyond the core set, a few basic functions appear in Office-suite questions.
\texttt{COUNTIF} counts the cells that meet one condition, and \texttt{SUMIF}
adds the cells that meet one condition; both extend COUNT and SUM. On the
statistics side, \texttt{MEDIAN} returns the middle value of the numbers,
while \texttt{MAX} and \texttt{MIN} return the ends of the range.
\texttt{ROUND} rounds a number to a stated number of digits. \texttt{LEN}
counts the characters in a text value. \texttt{TODAY} returns the current
date, and \texttt{NOW} returns the current date and time.

\begin{center}
\begin{tabular}{p{0.22\textwidth}p{0.62\textwidth}}
\toprule
\textbf{Function} & \textbf{What it returns} \\
\midrule
\texttt{COUNTIF} & The number of cells meeting one condition. \\
\texttt{SUMIF} & The total of the cells meeting one condition. \\
\texttt{MEDIAN} & The middle value of the numbers. \\
\texttt{ROUND} & A number rounded to a set number of digits. \\
\texttt{LEN} & The number of characters in a text value. \\
\texttt{TODAY} & The current date. \\
\bottomrule
\end{tabular}
\end{center}

At JKSSB depth, the exam is satisfied by knowing what each function returns
and that every one of them is written inside a formula that starts with
\texttt{=}.

\section{AutoSum}
\textbf{AutoSum} is a button, not a function, that quickly inserts a
\texttt{SUM} formula. When the active cell is placed below or beside a column
or row of numbers, AutoSum looks at the adjacent numbers and writes a
\texttt{=SUM(...)} formula over the range it selects. Pressing the button
therefore saves the user from typing the range by hand. The keyboard shortcut
for AutoSum is \texttt{Alt + =}. Its drop-down arrow also offers
\texttt{Average}, \texttt{Count Numbers}, \texttt{Max}, and \texttt{Min}, so
the same button can insert any of the core functions, not only SUM.

\begin{definitionbox}{AutoSum}
An Excel button that automatically inserts a \texttt{SUM} formula for the
range next to the active cell; the shortcut is \texttt{Alt + =}.
\end{definitionbox}

\section{Combining functions in one formula}
Functions can be nested inside one another, so the result of one function
becomes the argument of another. For example,
\texttt{=ROUND(AVERAGE(A1:A10),1)} first finds the mean of A1:A10 and then
rounds it to one decimal place. In the same way,
\texttt{=IF(COUNT(B1:B5)=0,0,AVERAGE(B1:B5))} shows 0 when the range contains
no numbers at all, instead of producing an error, and otherwise shows the mean
of the numbers present. Nesting is read from the inside out: the innermost
function is calculated first, and its result is passed outwards to the next
function.

\begin{example}
Consider \texttt{=IF(COUNT(C1:C6)=0,"No data",AVERAGE(C1:C6))}. If C1:C6 holds
only text or blanks, \texttt{COUNT} returns 0, the test is true, and the cell
shows \texttt{No data}. If the range holds numbers, the test is false and the
cell shows their average. This pattern is common on real marks sheets, where a
column may be empty early in the term.
\end{example}

The arguments of a function can also be single cells, several separate cells,
whole ranges, or even constant values such as \texttt{=SUM(A1, 5, B3:B7)}. A
reference such as \texttt{A1} can be relative or absolute (\texttt{\$A\$1});
that choice does not change what the function does, only how the reference
behaves when the formula is copied.

\section{Error values}
When a formula cannot be calculated, Excel displays an \emph{error value}
beginning with the hash sign \texttt{\#}. Each error points to a different
kind of problem, and recognising them is a common direct-recall question.

\begin{center}
\begin{tabular}{p{0.22\textwidth}p{0.60\textwidth}}
\toprule
\textbf{Error value} & \textbf{What it means} \\
\midrule
\texttt{\#DIV/0!} & A number is divided by zero or by a blank cell. \\
\texttt{\#VALUE!} & The wrong type of argument is used --- for example, text where a number is needed. \\
\texttt{\#NAME?} & Excel does not recognise a name in the formula, often a misspelt function name. \\
\texttt{\#REF!} & A cell reference is not valid, usually because the referenced cell was deleted. \\
\texttt{\#N/A} & The value is not available to the formula, as when a lookup finds no match. \\
\texttt{\#NUM!} & A number in the formula is not valid for the operation. \\
\texttt{\#NULL!} & A space (intersection) between two ranges that do not intersect. \\
\bottomrule
\end{tabular}
\end{center}

Two clarifications prevent confusion. First, the string \texttt{\#\#\#\#\#}
that sometimes appears in a cell is \emph{not} an error value; it only means
the column is too narrow to display the number, and it disappears when the
column is widened. Second, \texttt{\#DIV/0!} also appears when the divisor is
a blank cell, because a blank is read as zero in the division. Reading each
error as a clue to what Excel could not do is the fastest way to answer these
items.

\begin{example}
The error value points at the mistake that produced it. A formula
\texttt{=A1/0}, or \texttt{=A1/B1} where B1 is blank, shows
\texttt{\#DIV/0!}. A formula \texttt{=A1+B1} where A1 holds text shows
\texttt{\#VALUE!}. A misspelt function such as \texttt{=SUMM(A1:A5)} shows
\texttt{\#NAME?}. Deleting a cell that a formula refers to turns the reference
into \texttt{\#REF!}. A lookup that finds no match leaves the result
\texttt{\#N/A}, and a mathematically invalid argument such as
\texttt{=SQRT(-1)} gives \texttt{\#NUM!}.
\end{example}

\section{Exam retrieval and common confusions}
Five distinctions prevent most errors on this topic.
\begin{enumerate}
  \item A \textbf{formula} starts with \texttt{=}; a \textbf{function} is a
    named built-in operation used inside a formula. Without the equal sign,
    the entry is text (TRAP-30).
  \item \textbf{COUNT} counts numbers only; \textbf{COUNTA} counts all
    non-empty cells. Do not swap this near-neighbour pair.
  \item \textbf{AVERAGE} divides by the number of numeric cells, not by the
    total number of cells.
  \item A \textbf{blank} cell is read as zero in a numeric test, so a blank
    fails a ``greater than'' condition; text used in arithmetic gives
    \texttt{\#VALUE!} (KB PYQ-14).
  \item \textbf{AutoSum} inserts \texttt{=SUM(...)} and its shortcut is
    \texttt{Alt + =}, not a Ctrl combination.
\end{enumerate}

\begin{example}
An item asks which function counts a range that mixes numbers and labels.
Read the function name carefully: \texttt{COUNT} returns only the numeric
cells, while \texttt{COUNTA} returns every non-empty cell. If the stem asks
for the count of entries of any kind, the answer is \texttt{COUNTA}; if it
asks for the count of numbers, the answer is \texttt{COUNT}. A second item
may ask which error appears when a cell is divided by a blank; the blank
counts as zero, so the error is \texttt{\#DIV/0!}.
\end{example}

\subsection*{Exercises}
\begin{enumerate}[label=\arabic*.]
  \item Explain the difference between a formula and a function in MS Excel.
  \item A range A1:A8 contains 5 numbers, 2 text labels, and 1 blank cell.
    State the values returned by \texttt{=COUNT(A1:A8)} and
    \texttt{=COUNTA(A1:A8)}.
  \item Write the general syntax of the IF function and explain each of its
    three arguments.
  \item A cell holds \texttt{=IF(A1>10,"Yes",IF(A1>0,"Small","No"))} and A1 is
    blank. What is shown, and why?
  \item Match each error value with its cause: \texttt{\#DIV/0!},
    \texttt{\#VALUE!}, \texttt{\#NAME?}, \texttt{\#REF!}.
\end{enumerate}

\subsection*{Solutions}
\begingroup\small
\begin{enumerate}[label=\arabic*.]
  \item A formula is any calculation that begins with \texttt{=}, such as
    \texttt{=A1+A2}. A function is a predefined, named operation such as
    \texttt{SUM} or \texttt{IF} that is used inside a formula. Every function
    is written as a formula, but a formula need not contain a function.
  \item \texttt{=COUNT(A1:A8)} returns 5, because only the numeric cells are
    counted. \texttt{=COUNTA(A1:A8)} returns 7, because the 5 numbers and the
    2 text labels are all non-empty; both ignore the blank.
  \item \texttt{=IF(logical\_test, value\_if\_true, value\_if\_false)}. The
    first argument is a test answered true or false; the second is the value
    returned when it is true; the third is the value returned when it is
    false.
  \item The formula shows \texttt{No}. A blank is read as zero in a numeric
    test, so A1 is neither greater than 10 nor greater than 0, both tests are
    false, and the final value \texttt{No} is returned.
  \item \texttt{\#DIV/0!} --- division by zero or by a blank cell;
    \texttt{\#VALUE!} --- the wrong type of argument, such as text where a
    number is needed; \texttt{\#NAME?} --- an unrecognised name, often a
    misspelt function; \texttt{\#REF!} --- an invalid cell reference after
    the referred cell was deleted.
\end{enumerate}
\endgroup
\end{document}
